\scanpage{8}
\mspage{7}
On vient de voir que $a)\unsure{\Rightarrow|=} b)$ prouvons l'inverse, supposons $b)$ prouvons $a)$ i.e. soient $A\in\Psi$, $B\in\unsure{\mathscr{L}|\mathcal{L}}$, $B\deleted{\unsure{\in}}\subset A$ prouvons $B\in\Psi$

Considérons $B^\circ$, on a $B^\circ\ne\varnothing$ ($\Psi$ adm.) et comme $R_\Psi=R_{\mathscr{L}}\lvert\deleted{\unsure{\mathcal{P}|B}}\Psi\rvert$ \inserted{\unsure{saturé|saturée}}, on a $B^\circ=C^\circ$ pour $C\in\deleted{\mathscr{L}}\Psi$ convenable, mais cela implique $\deleted{\unsure{C}}B=C$, donc $B\in\Psi$, qed.

On dit que $\Psi$ est une \uline{sous-famille} \uline{admissible}.

\uline{Scholie} Elles correspondent \inserted{(avec conservation des rel. d'ordre)} aux parties fermées de $\mathscr{L}$, ou de $|\mathscr{L}|$. \unsure{\ensuremath{\Rightarrow}|=}\,Elles sont \uline{stables} par $\cup$ et $\cap$ quelc. dans $|\mathscr{L}|$.

\uline{Proposition} Soient $\mathscr{L},\Psi\subset\unsure{\mathcal{P}|\mathscr{P}}(L)$ deux ens de parties admissibles. Conditions équiv.

a) $\mathscr{L}\cup\Psi$ est admissible, et $\mathscr{L},\Psi$ en sont \unsure{deux|des} sous-ens. admissibles (i.e. fermés)

\inserted{$\deleted{\unsure{a}}a')$ $\mathscr{L}$ et $\Psi$ sont \deleted{\unsure{deux}}sous-familles adm. d'une famille \inserted{adm.} $\Sigma$}
b) $\mathscr{L}\cap\Psi$ est $\mathscr{L}$-fermé dans $\unsure{\Phi|\mathscr{L}}$ et \inserted{$\Psi$-fermé} dans $\Psi$,
\[
 \text{et}\qquad |\mathscr{L}|\cap|\Psi|=|\mathscr{L}\cap\Psi|
\]
c) $|\mathscr{L}|\cap|\Psi|$ est $\mathscr{L}$-fermé dans $|\mathscr{L}|$ et \unsure{\ensuremath{\{}|,} $\Psi$-fermé dans $\Psi$, et \deleted{\unsure{toute} $A\in\mathscr{L}$ telle que} pour $A\subset|\mathscr{L}|\cap|\Psi|$, $A\in\mathscr{L}$ \deleted{de $\mathscr{L}$} ssi $\unsure{A|1}\in\Psi$, (i.e. ssi $A\in\mathscr{L}\cap\Psi$).

% slanted
\margin{Cf condition équiv. page 11 (\unsure{genre|sorte} d), ``strate par strate''}
\deleted{d) $\forall A\in\mathscr{L}, B\in\Psi$, on ait $\unsure{A}\ \mathscr{L}_{\unsure{\leq A|<A}}$ et $\Psi_{\unsure{\leq|<}}$}

$a')\Leftrightarrow a)$ évident par les sorites précédents de \unsure{\ensuremath{=}|m} que $a)$ implique $b)$ et $c)$. Reste à prouver $b\Rightarrow c$, $c\Rightarrow a$.

Prouvons d'abord $b\Rightarrow c$. Comme $|\mathscr{L}|\cap|\Psi|=\unsure{\lvert\mathscr{L}\cap\Psi)|\lvert\mathscr{L}\cap\Psi\rvert}$ par $b)$, et $\mathscr{L}\cap\Psi$ fermé dans $\mathscr{L}$ et
